% SIGN_NEXT main-critical contribution. Author: GPT-6 Astra Pro, child /root/critical_lefschetz.
% Reference keys to integrate: GrothCoates1966, GrothCoates1967, Kleiman1968, SGA5.
\section{The critical-degree proposition: a corrected formulation}
\label{sec:critical-proposition}

Throughout this section, $H$ is a characteristic-zero Weil cohomology
with hard Lefschetz and the usual weak Lefschetz theorem for the smooth
hyperplane sections in use.  Tate twists are suppressed in formulas; an
algebraic class of degree $2r$ belongs to $H^{2r}(T)(r)$.
Write
\[
 \operatorname{Alg}^r(T)=\operatorname{im}\bigl(
 \mathrm{CH}^r(T)_{\mathbf Q}\longrightarrow H^{2r}(T)(r)\bigr).
\]
Algebraicity always refers to rational cycle classes.  We use
$C_{\mathrm{Gr}}(X)$ for the condition called $C(X)$ in the Coates
letters: the normalized Lefschetz lowering correspondence is algebraic.
It is the condition called $B(X)$ by Kleiman in 1968 and in modern
standard-conjecture notation.  It is not the modern K\"unneth condition.
The cycle condition $A(X,L)$ says that the hard-Lefschetz isomorphisms
restrict to isomorphisms on the spaces $\operatorname{Alg}^r(X)$.

\paragraph{SN-C01 (the source and the left-hand label).}
The left side of the proposition relevant to the task is
$C_{\mathrm{Gr}}(X)$, not the weaker cycle condition $A(X,L)$.
This is the mathematical reconstruction of the ambiguous overtyped
letter in cote~16, scan~050 (folio~49).  It is corroborated by the
CSG transcription of the letter dated 6.1.1966, printed p.~3,
\cite{GrothCoates1966}, and by the preceding archived letter dated 4.1.1967,
printed p.~2, Proposition (i), (ii bis), \cite{GrothCoates1967}.
The exact published result to use is Kleiman's Theorem~2.13(i),
(ii$'$), pp.~372--373 \cite{Kleiman1968}: if $Y\subset X$ is a
smooth hyperplane section and $\dim X=n$, then
\begin{equation}
 C_{\mathrm{Gr}}(X)\quad\Longleftrightarrow\quad
 C_{\mathrm{Gr}}(Y)\ \text{and}\
 L_{X\times X}^{\,2}:
 \operatorname{Alg}^{n-1}(X\times X)
 \xrightarrow{\ \sim\ }
 \operatorname{Alg}^{n+1}(X\times X),
 \label{eq:kleiman-critical}
\end{equation}
where the product is polarized by the Segre class
$L_{X\times X}=L_X\otimes1+1\otimes L_X$.
This known theorem is cited, not reproved here.  In particular,
a statement about $A(X,L)$ alone must not be used in its place.
The same source's condition (iv) gives the criterion along a flag
of smooth hyperplane sections.

There is a second issue in scan~050, distinct from the left-hand letter.
Its inequalities use $2i$ while several cohomological superscripts
are printed $i$; the critical-degree prescription says $2i=d-2$
when $d$ is even, whereas the preceding ordinary-Gysin variant uses
the upper half of the range.  One cannot silently combine these
pieces into a valid ordinary-Gysin condition.  The following
counterexample requires $i=0$, so it is unaffected by the missing
factor of two in the superscripts.

\begin{proposition}[SN-C02: the literal lower critical Gysin condition fails]
\label{prop:literal-critical-fails}
Read the even critical ordinary-Gysin condition as
\begin{equation}
 j_*\operatorname{Alg}^{m-1}(W)=\operatorname{Alg}^{m}(T),
 \qquad \dim T=2m,
 \label{eq:bad-low-gysin}
\end{equation}
for a smooth hyperplane section $j:W\hookrightarrow T$.
It is not a consequence of $C_{\mathrm{Gr}}(X)$ when $T=X\times X$,
and therefore cannot occur on the right side of an equivalence with
that condition.  This failure occurs for $X=\mathbf P^1$ and also
for an elliptic curve.
\end{proposition}
\begin{proof}
Take $X=\mathbf P^1$, so that $T=\mathbf P^1\times\mathbf P^1$,
$m=1$, and let $a,b$ be the two ruling classes.  Take $W$ to be a
smooth divisor of bidegree $(1,1)$.  It is isomorphic to
$\mathbf P^1$, and
\[
 H^*(T,\mathbf Q)=\mathbf Q[a,b]/(a^2,b^2),\qquad
 j_*1=a+b.
\]
Every cohomology class in question is algebraic.  The left side of
\eqref{eq:bad-low-gysin} is the line $\mathbf Q(a+b)$, while its
right side is $\mathbf Qa\oplus\mathbf Qb$.  They are not equal.
The restriction in degree zero is an isomorphism and satisfies the
critical $A'$ condition.  Both $X$ and its zero-dimensional hyperplane
section satisfy $C_{\mathrm{Gr}}$: on $\mathbf P^1$, the fundamental
class of $\mathbf P^1\times\mathbf P^1$ is the correspondence taking
the degree-two generator to $1$ and vanishing on the other degrees.
Moreover $A(T)^0$ itself holds, since
$(a+b)^2\cdot1=2ab$.  Thus the failed weak condition cannot be
attributed to a failure of the genuine critical Lefschetz statement.

For an elliptic curve $E$, choose any sufficiently ample product
polarization on $E\times E$ and a smooth connected member $W$.
The two factor divisor classes are independent: they have squares
zero and intersection one.  Hence
$\dim_{\mathbf Q}\operatorname{Alg}^1(E\times E)\ge2$, whereas
$j_*H^0(W)$ has dimension one.  The same failure follows.  For a
curve, $C_{\mathrm{Gr}}$ holds directly: the only nonzero lowering
map is $H^2\to H^0$, induced, after division by the degree of the
polarization, by the fundamental product correspondence.
\end{proof}

We now give a replacement stated in terms of actual maps.  It uses
the final even-dimensional diagram in scan~050, but it makes its
extra $L_T$ part of the hypothesis.  This is a mathematical correction,
not a claim that the preceding definition of $A''$ in the manuscript
already contained that factor.

\begin{theorem}[SN-C03: the exact even critical factorization]
\label{thm:weighted-critical}
Let $T$ be smooth projective of dimension $2m\ge2$, and let
$j:W\hookrightarrow T$ be a smooth hyperplane section with class $L_T$.
Put $L_W=j^*L_T$.  Define the following conditions:
\begin{align*}
 R_0(T,W):\ &x\in H^{2m-2}(T),\quad
             j^*x\in\operatorname{Alg}^{m-1}(W)
             \ \Longrightarrow\ x\in\operatorname{Alg}^{m-1}(T);\\
 G_{L,0}(T,W):\ &L_Tj_*\operatorname{Alg}^{m-1}(W)
                    =\operatorname{Alg}^{m+1}(T);\\
 A_0(T):\ &L_T^2:\operatorname{Alg}^{m-1}(T)
                    \xrightarrow{\ \sim\ }
                    \operatorname{Alg}^{m+1}(T).
\end{align*}
Then
\begin{equation}
 A_0(T)\quad\Longleftrightarrow\quad
 R_0(T,W)\ \text{and}\ G_{L,0}(T,W).
 \label{eq:weighted-equivalence}
\end{equation}
No algebraicity of the inverse Lefschetz operator on $W$ is an
additional hypothesis in this equivalence.

For comparison, put
\[
 G_0(T,W):\quad
 j_*\operatorname{Alg}^{m}(W)=\operatorname{Alg}^{m+1}(T),
 \qquad
 A_0(W):\quad L_W:\operatorname{Alg}^{m-1}(W)
                       \xrightarrow{\ \sim\ }\operatorname{Alg}^{m}(W).
\]
There is an exact further equivalence
\begin{equation}
 G_{L,0}(T,W)\quad\Longleftrightarrow\quad
 G_0(T,W)\ \text{and}\ A_0(W).
 \label{eq:gysin-hidden-inverse}
\end{equation}
Thus replacing $G_{L,0}$ by ordinary upper critical Gysin
surjectivity drops precisely the indicated inverse-Lefschetz
assertion on the hyperplane $W$.
\end{theorem}
\begin{proof}
Weak Lefschetz and its Poincar\'e-dual form give isomorphisms
\[
 r=j^*:H^{2m-2}(T)\xrightarrow{\sim}H^{2m-2}(W),\qquad
 g=j_*:H^{2m}(W)\xrightarrow{\sim}H^{2m+2}(T).
\]
The projection and self-intersection formulas give
\begin{equation}
 L_T^2=gL_Wr=L_Tj_*r.
 \label{eq:critical-factorization}
\end{equation}
All the displayed maps preserve algebraic cycle classes whenever
source and target degrees match.

Suppose first that $A_0(T)$ holds.  If $r(x)$ is algebraic, then
$L_Tj_*r(x)=L_T^2x$ is algebraic.  Surjectivity of $A_0(T)$ supplies
an algebraic $x_0$ with $L_T^2x_0=L_T^2x$.  Hard Lefschetz makes
$L_T^2$ injective on all cohomology, so $x=x_0$.  This proves
$R_0(T,W)$.  If $\beta\in\operatorname{Alg}^{m+1}(T)$, write
$\beta=L_T^2x_0$ with $x_0$ algebraic.  Then $r(x_0)$ is algebraic
and \eqref{eq:critical-factorization} proves $G_{L,0}(T,W)$.

Conversely, assume $R_0$ and $G_{L,0}$.  Given an algebraic
$\beta$ in the upper critical degree, choose an algebraic
$y\in H^{2m-2}(W)$ with $L_Tj_*y=\beta$.  The isomorphism $r$
gives a unique $x$ with $r(x)=y$.  Condition $R_0$ makes $x$
algebraic, and \eqref{eq:critical-factorization} gives
$L_T^2x=\beta$.  Injectivity is again hard Lefschetz.  This proves
\eqref{eq:weighted-equivalence}.

Next assume $G_{L,0}$.  Since $L_Tj_*=j_*L_W$, it immediately
implies $G_0$.  For any algebraic $y\in H^{2m}(W)$, the class
$g(y)$ is algebraic.  Apply $G_{L,0}$ to choose algebraic
$z\in H^{2m-2}(W)$ with $g(y)=gL_Wz$.  Injectivity of $g$ gives
$y=L_Wz$.  Hard Lefschetz on $W$ gives the injectivity needed for
$A_0(W)$.  In the other direction, $G_0$ supplies, for an algebraic
$\beta$ on $T$, an algebraic $y$ with $g(y)=\beta$; $A_0(W)$
supplies an algebraic $z$ with $L_Wz=y$.  Therefore
$\beta=gL_Wz=L_Tj_*z$.  This proves
\eqref{eq:gysin-hidden-inverse}.
\end{proof}

\begin{corollary}[SN-C04: a usable corrected proposition for scan 050]
\label{cor:corrected-scan050}
Let $X$ be smooth projective of dimension $n\ge1$, let $Y\subset X$
be a smooth hyperplane section, and polarize $T=X\times X$ by
$L_T=L_X\otimes1+1\otimes L_X$.  Let $W\subset T$ be a smooth
hyperplane section for this polarization.  Assume the hard and weak
Lefschetz hypotheses of Kleiman's Theorem~2.13.  Then
\begin{equation}
 \boxed{\quad C_{\mathrm{Gr}}(X)\ \Longleftrightarrow\
 C_{\mathrm{Gr}}(Y)+R_0(T,W)+G_{L,0}(T,W).\quad}
 \label{eq:corrected-scan050}
\end{equation}
Along a flag of smooth hyperplane sections of $X$, the conjunction
of the two corrected weak conditions on each square is equivalent
to $C_{\mathrm{Gr}}(X)$.  The hyperplane $W\subset X\times X$
is not $Y\times Y$; their dimensions are $2n-1$ and $2n-2$,
respectively.
\end{corollary}
\begin{proof}
Apply Theorem~\ref{thm:weighted-critical} to $T$, with $m=n$.
It identifies the two displayed weak conditions with exactly the
critical condition in \eqref{eq:kleiman-critical}.  Kleiman's
Theorem~2.13(i), (ii$'$) now proves the assertion.  Iterating down
the flag and using the same theorem's (iv), or the fact that
$C_{\mathrm{Gr}}$ holds in dimension at most one, proves the final
assertion.  No implication from the uncorrected ordinary-Gysin
condition is used.
\end{proof}

\begin{proposition}[SN-C05: the upper-degree repair alone does not justify the diagram argument]
\label{prop:formal-critical-obstruction}
Restriction detection $R_0$ and ordinary upper critical Gysin
surjectivity $G_0$, together with the cohomological hard-Lefschetz
isomorphisms and the restriction--Gysin identities, do not formally
imply $A_0$.  This is a statement about those linear-algebra data,
not a geometric counterexample to a standard conjecture.
\end{proposition}
\begin{proof}
Take a Lefschetz module $V$ with vectors $v_0,v_1,v_2$ in degrees
$2m-2,2m,2m+2$ and $L_Vv_0=v_1$, $L_Vv_1=v_2$, $L_Vv_2=0$.
Take $U$ with vectors $u_0,u_1$ in degrees $2m-2,2m$, and
$L_Uu_0=u_1$, $L_Uu_1=0$.  Define
\[
 rv_0=u_0,\quad rv_1=u_1,\quad rv_2=0,
 \qquad gu_0=v_1,\quad gu_1=v_2.
\]
Then $gr=L_V$, $rg=L_U$, $rL_V=L_Ur$ and $gL_U=L_Vg$.
Both modules satisfy hard Lefschetz about their indicated centers,
and the low restriction and high Gysin maps are isomorphisms.
Declare their ``algebraic'' subspaces to be
\[
 V_{\mathrm{alg}}^{2m-2}=0,\quad
 V_{\mathrm{alg}}^{2m}=\mathbf Qv_1,\quad
 V_{\mathrm{alg}}^{2m+2}=\mathbf Qv_2,
 \qquad
 U_{\mathrm{alg}}^{2m-2}=0,\quad
 U_{\mathrm{alg}}^{2m}=\mathbf Qu_1.
\]
Every defined map preserves these subspaces.  Restriction detection
holds in the low degree because $r$ is injective and the target
algebraic subspace is zero.  The high Gysin map carries
$\mathbf Qu_1$ onto $\mathbf Qv_2$, so $G_0$ holds.  Nevertheless
$L_V^2$ takes the zero algebraic subspace in the low degree to zero,
not to $\mathbf Qv_2$; hence $A_0$ fails.  Equivalently,
$L_U^{-1}(u_1)=u_0$ is the missing algebraic preimage.
\end{proof}

\paragraph{Placement in Grothendieck's frameworks.}
The replacement \eqref{eq:corrected-scan050} is a statement about
rational algebraic correspondences and cycle classes, at precisely
cote~16 scans~034, 036, 040, 049--050.  Its restriction and Gysin
maps are the maps in weak Lefschetz (SGA~5, Expos\'e~VII,
Theorem~7.1, and the dual form), not information recovered from
Frobenius characteristic polynomials.  For a smooth model over an
open subscheme of $\operatorname{Spec}\mathbf Z$, the same maps
are restriction and proper pushforward on the corresponding
\(\ell\)-adic realizations, with the usual twists.  The
$\mathbf P^1\times\mathbf P^1$ counterexample and its smooth
$(1,1)$ diagonal are already defined over $\mathbf Z$, so the
rank failure occurs in every geometric fibre.  In rational
$K$-theory, the extra $L_T$ is multiplication by the first-Chern-class
operation in the associated graded for the gamma filtration;
forgetting it changes codimension by one.  In the category of
homological motives, $C_{\mathrm{Gr}}(X)$ supplies the inverse
Lefschetz morphisms.  None of SN-C01--SN-C05 alone asserts the
index inequality, positivity of an involution, or semisimplicity;
those are separate steps on scan~049.  Nor do these replacements
change the already-established sign-only purity theorem or any
local $L$-factor.
